% TOP_PROBLEM: {"problemId": "problem.riemann-hypothesis", "canonicalTitle": "Riemann Hypothesis", "releaseRank": 2, "primaryDomain": "number_theory_arithmetic_geometry", "primaryDomainLabel": "Number theory & arithmetic geometry", "displayStatus": "open", "releaseStatus": "open"}
% !TeX program = lualatex
% Definitions and sources retained from research_batch01.tex; research is in GPT_6_Astra_Ultra.
\documentclass[11pt]{article}
\usepackage[margin=25mm]{geometry}
\usepackage{fontspec}
\setmainfont{Latin Modern Roman}
\usepackage{amsmath,amssymb,mathtools}
\usepackage{unicode-math}
\setmathfont{Latin Modern Math}
\usepackage{enumitem,needspace}
\usepackage{xurl,hyperref}
\hypersetup{colorlinks=true,urlcolor=blue}
\setcounter{secnumdepth}{1}
\setlength{\parindent}{0pt}
\setlength{\parskip}{5pt}
\setlength{\emergencystretch}{3em}
\setlist[itemize]{leftmargin=3em,itemsep=5pt,topsep=4pt,parsep=0pt}
\allowdisplaybreaks
\newcommand{\N}{\mathbb{N}}
\newcommand{\Z}{\mathbb{Z}}
\newcommand{\Q}{\mathbb{Q}}
\newcommand{\R}{\mathbb{R}}
\newcommand{\C}{\mathbb{C}}
\newcommand{\F}{\mathbb{F}}
\newcommand{\A}{\mathbb{A}}
\newcommand{\PP}{\mathbb P}
\newcommand{\E}{\mathbb{E}}
\newcommand{\eps}{\varepsilon}
\newcommand{\dd}{\,\mathrm d}
\newcommand{\sref}[2]{\hyperlink{source:#1:#2}{[S#2]}}
\newcommand{\eref}[2]{\hyperlink{extra:#1:#2}{[E#2]}}
% Turn the next switch off for a shorter reading copy. The quotations preserve
% every exactTarget field, including qualifications omitted from a short summary.
\newif\ifshowcatalogue
\showcataloguetrue
\newcommand{\cataloguescope}[1]{\ifshowcatalogue
 \paragraph{Exact scope in the supplied catalogue.}
 \begin{quote}\small #1\end{quote}\fi}
\begin{document}
\setcounter{section}{1}
% ================================================================
% problemId: problem.riemann-hypothesis
% source releaseRank: 2; source releaseStatus: open
% exactTarget SHA-256: 43a61dff12aa4263f3e68f91192ddf29e0218c93b508ad4fad99b54d1df28674
\clearpage
\section{\texorpdfstring{Riemann Hypothesis}{Riemann Hypothesis}}\label{2}
\subsection{Definitions and mathematical statement}
For $s\in\C$ with $\Re s>1$, define
\[
 \zeta(s)=\sum_{n=1}^{\infty}n^{-s}=\prod_{p\text{ prime}}(1-p^{-s})^{-1}.
\]
Use its meromorphic continuation to $\C$, whose only pole is at $s=1$. The assertion is
\[
 \forall\rho\in\C,\qquad \bigl(0<\Re\rho<1\ \text{and}\ \zeta(\rho)=0\bigr)
 \ \Longrightarrow\ \Re\rho=\tfrac12.
\]
The open critical strip excludes the negative even integer zeros. The assertion concerns locations, not multiplicities, of zeros.
\par\smallskip\textit{Formulation and concept references:} \sref{2}{1}.
\cataloguescope{For the meromorphic continuation of the Riemann zeta function zeta(s)=sum\_\{n>=1\} n\textasciicircum{}(-s), initially defined for Re(s)>1, every zero s in the open critical strip 0<Re(s)<1 has Re(s)=1/2. This is the Riemann Hypothesis for zeta alone; it does not assert simplicity of its zeros.}
\subsection{Short English statement}
Are all zeta zeros between the vertical lines zero and one located on the halfway line?
\subsection{Sources}
\begin{itemize}\raggedright
\item[\textnormal{[S1]}]\hypertarget{source:2:1}{} Enrico Bombieri, Problems of the Millennium: the Riemann Hypothesis, Section I and the prime-counting equivalence.
\url{https://www.claymath.org/wp-content/uploads/2022/05/riemann.pdf}.
{\small\textit{Catalogue formulation source. Consulted 24 September 2026: Section I, definition and zero-location statement.}}
\item[\textnormal{[S2]}]\hypertarget{source:2:2}{} Riemann Hypothesis.
\url{https://www.claymath.org/millennium/riemann-hypothesis/}.
{\small\textit{Catalogue research/status reference; not a proof endorsement.}}
\item[\textnormal{[S3]}]\hypertarget{source:2:3}{} A new proof that more than 2/3 of the zeros of the Riemann zeta function are simple and on the critical line.
\url{https://arxiv.org/abs/2609.02882}.
{\small\textit{Catalogue research/status reference; not a proof endorsement.}}
\item[\textnormal{[S4]}]\hypertarget{source:2:4}{} Lamzouri lecture at Max Planck Institute for Mathematics, September 3, 2026.
\url{https://math-events.uni-bonn.de/event/1560/}.
{\small\textit{Catalogue research/status reference; not a proof endorsement.}}
\item[\textnormal{[S5]}]\hypertarget{source:2:5}{} On the Extended, Generalized, and Grand Riemann Hypotheses: A Unified Approach Based on the General Properties of L-Functions — Weicun Zhang, v19.
\url{https://www.preprints.org/manuscript/202506.0481}.
{\small\textit{Catalogue research/status reference; not a proof endorsement.}}

% BEGIN ADDED RESEARCH SOURCES -- batch 1, problem 2
\item[\textnormal{[E1]}]\hypertarget{extra:2:1}{} Dave Platt and Tim Trudgian,
\emph{The Riemann hypothesis is true up to $3\cdot10^{12}$},
arXiv:2004.09765v1 (21 April 2020), Theorem~1 and Section~2;
published in \emph{Bulletin of the London Mathematical Society} (2021).
\url{https://arxiv.org/pdf/2004.09765v1}.
{\small\textit{Research reference: rigorous bounded-height verification,
including completeness of the zero count. Consulted 24 September 2026.}}

\item[\textnormal{[E2]}]\hypertarget{extra:2:2}{} Larry Guth and James Maynard,
\emph{New large value estimates for Dirichlet polynomials},
\emph{Annals of Mathematics} 203 (2026), no.~2, 623--675,
DOI: 10.4007/annals.2026.203.2.6, Theorem~1.2.
\url{https://annals.math.princeton.edu/2026/203-2/p06}.
Author manuscript: arXiv:2405.20552v2 (7 April 2026),
\url{https://arxiv.org/pdf/2405.20552v2}.
{\small\textit{Research reference: zero-density estimate; publication metadata
and the stated theorem were checked. Consulted 24 September 2026.}}

\item[\textnormal{[E3]}]\hypertarget{extra:2:3}{} Michael Griffin, Ken Ono,
Larry Rolen and Don Zagier,
\emph{Jensen polynomials for the Riemann zeta function and other sequences},
arXiv:1902.07321v2 (31 March 2019), Section~1, Theorems~1--2;
published in \emph{Proceedings of the National Academy of Sciences} (2019).
\url{https://arxiv.org/pdf/1902.07321v2}.
{\small\textit{Research reference: the fixed-degree, sufficiently-large-shift
quantifiers, not all-degree hyperbolicity. Consulted 24 September 2026.}}

\item[\textnormal{[E4]}]\hypertarget{extra:2:4}{} Brad Rodgers and Terence Tao,
\emph{The de Bruijn--Newman constant is non-negative},
\emph{Forum of Mathematics, Pi} 8 (2020), e6,
DOI: 10.1017/fmp.2020.6.
Author manuscript arXiv:1801.05914v5 (3 July 2021),
Section~1, equations~(1)--(4), and the main theorem.
\url{https://arxiv.org/pdf/1801.05914v5}.
{\small\textit{Research reference: the theta-kernel normalization and the
endpoint obstruction. Consulted 24 September 2026.}}

\item[\textnormal{[E5]}]\hypertarget{extra:2:5}{} David W. Farmer,
\emph{Jensen polynomials are not a plausible route to proving the Riemann
Hypothesis}, arXiv:2008.07206, Sections~2--4.
\url{https://arxiv.org/pdf/2008.07206}.
{\small\textit{Research reference: differentiation and information loss;
the paper's methodological assessment is not used as an impossibility theorem.
Consulted 24 September 2026.}}

\item[\textnormal{[E6]}]\hypertarget{extra:2:6}{} Jonathan Holland,
\emph{A new hyperbolicity wedge and a joint semicircle limit for Jensen
polynomials of Riemann's $\xi$-function},
arXiv:2608.08682v1 (9 August 2026), Theorem~1.1.
\url{https://arxiv.org/html/2608.08682v1}.
{\small\textit{Research reference: recent manuscript, statement inspected,
full proof not independently audited here; not used as a premise of any lemma.
Consulted 24 September 2026.}}

\item[\textnormal{[E7]}]\hypertarget{extra:2:7}{} Kiran S. Kedlaya,
\emph{More on the zeroes of $\zeta$}, MIT 18.785: Analytic Number Theory,
spring 2007, Section~1, Lemma~1 and Theorems~2--3.
\url{https://kskedlaya.org/18.785/zeroes.pdf}.
{\small\textit{Research reference: growth of the completed function, zero
counting, and Hadamard factorization. Consulted 24 September 2026.}}

\item[\textnormal{[E8]}]\hypertarget{extra:2:8}{} Paul Garrett,
\emph{12. Weierstrass and Hadamard products}, author lecture notes,
3 February 2021, Section~3 (Hadamard products).
\url{https://www-users.cse.umn.edu/~garrett/m/complex/notes_2020-21/12_Hadamard_products.pdf}.
{\small\textit{Research reference: canonical products and the finite-order
factorization theorem. Consulted 24 September 2026.}}

\item[\textnormal{[E9]}]\hypertarget{extra:2:9}{} Youness Lamzouri,
\emph{A new proof that more than $2/3$ of the zeros of the Riemann zeta
function are simple and on the critical line},
arXiv:2609.02882v2 (8 September 2026), Theorem~1.1,
Proposition~2.1 and Sections~2--3.
\url{https://arxiv.org/html/2609.02882v2}.
{\small\textit{Version-specific research supplement to [S3]: the stated bound
and method were inspected. Neither a complete analytic proof audit nor a
formal-certificate audit was performed. Consulted 24 September 2026.}}
% END ADDED RESEARCH SOURCES -- batch 1, problem 2
\end{itemize}

\end{document}
